\section{Microservice implementation: Indexing}
\label{sec:indexing}

Indexing turns the complete books of the datalake into the inverted index and the metadata in MongoDB, and decides
when a new version of the index becomes visible to Search. It is written in Java 17 with Javalin and Gson. Its HTTP
API, \code{api/indexing-openapi.yaml}, has five endpoints: \code{POST /index/update/\{book\_id\}} indexes one book,
\code{POST /index/publish} and \code{POST /index/rebuild} publish and rebuild the index, and \code{GET /index/status}
and \code{GET /status} report its state. Its MongoDB side is specified in \code{index-publication.md} and its
Hazelcast side in \code{coordination.md}, both in \code{api/contracts/}. Its main classes are:

\begin{table}[H]
\centering\small
\begin{tabularx}{\linewidth}{@{}L{5.3cm}>{\raggedright\arraybackslash}X@{}}
\toprule
\textbf{Class} & \textbf{Responsibility} \\
\midrule
\code{IndexingApi} & Validates the requests and turns results and exceptions into the HTTP answers of the contract \\
\code{IndexingService} & Reads, parses and tokenizes books, and writes them: one book (\code{update}) or a batch
  (\code{indexBatch}) \\
\code{IndexingWriter} & The consumer of \code{index-queue}: holds \code{index-writer}, claims batches, completes or
  fails each book \\
\code{LeaseRenewer} & Renews the leases of the batch being written \\
\code{Publisher} & Publishes a new version of the index with a conditional switch of \code{CURRENT} \\
\code{HazelcastPublicationNotifier} & After a publication, announces \code{IndexPublished} and marks the books
  \code{PUBLISHED}, in a thread of its own \\
\code{Rebuilder} & \code{POST /index/rebuild}: rebuilds the whole index into a new collection and switches to it \\
\code{MongoPublicationRepository} & The publication collections: \code{index\_meta}, \code{index\_books},
  \code{index\_publications}, \code{index\_log} \\
\bottomrule
\end{tabularx}
\caption{Structure of the Indexing Service. The datalake, the tokenizer, the metadata parser and the MongoDB stores of
the index and the metadata are in \code{common}, shared with the other services.}
\end{table}

\textbf{One book, one batch.} Both ways in index the same way, in two steps. \emph{Prepare} reads the header and the
body of each book from the datalake, parses the metadata (title, author, language, year) and tokenizes the body with the
tokenizer of Stage~1, the same one Search uses for queries: ASCII letters in lower case and digits, tokens of at least
2 characters, without stopwords. It runs in parallel on a pool of
\code{INDEXING\_THREADS} (4). \emph{Write} runs in one thread and in this order: the metadata of every book of the batch
with one bulk write into \code{books}, the posting lists with one bulk write into the inverted index, and \emph{last}
one entry per book in \code{index\_log}. The inverted index sends one \code{\$addToSet} upsert per \emph{distinct} term
of the batch, so a common word is written once per batch, not once per book. A book that cannot be prepared is left
out and reported, and never stops the others. A write that fails fails the whole batch.

\textbf{Why \code{index\_log} goes last.} It is the evidence that a book is complete in the index: the publication
takes its books from it, the book becomes \code{INDEXED} in Hazelcast only after it, and the recovery of the Control
(Section~\ref{sec:recovery}) reads it. A process that dies in the middle of a batch leaves no entry, so the book is
indexed again and never published half written.

\subsection{Readiness, task ownership and retries}

\textbf{Readiness.} Indexing reads a book only if its \code{\_READY} marker is in the datalake
(Section~\ref{sec:readiness}). Otherwise \code{POST /index/update} answers 404 and the queue writer fails the book with
\code{NOT\_READY}, which puts it back in the queue.

\textbf{Task ownership.} One thread per replica runs \code{IndexingWriter}. It first takes the Hazelcast lock
\code{index-writer}: only its holder consumes \code{index-queue}, and another replica waits as a standby and takes over
when the holder stops or dies. The holder claims a book, waiting for one, and then the books already waiting, without
waiting, up to \code{INDEX\_WRITE\_BATCH} (100). With few books a batch has one book and nothing waits for more; with a
backlog the batches fill up by themselves. After the write, it completes each indexed book (\code{INDEXED}). Each
claim has a lease of \code{LEASE\_SECONDS} (120~s). A batch on a large index can take longer than that (about 33~s per
batch at 10,000 books, Section~\ref{sec:benchmarks-indexing}), so \code{LeaseRenewer} renews the leases of the batch
every third of the lease while it is written. Without it, the Control would requeue books that are still being written.

\textbf{Retries.} Each failure goes back to Hazelcast with a code of the coordination contract:

\begin{table}[H]
\centering\small
\begin{tabularx}{\linewidth}{@{}L{4.6cm}L{4cm}X@{}}
\toprule
\textbf{Failure} & \textbf{Code} & \textbf{What happens} \\
\midrule
Book not complete in the datalake & \code{NOT\_READY} & Requeued, one attempt more \\
Book without indexable terms & \code{EMPTY\_BOOK} & \code{FAILED} at once; nothing is written \\
MongoDB does not answer & \code{STORAGE\_UNAVAILABLE} & The whole batch is requeued \emph{without} counting an attempt,
  and the writer waits 1, 2, 4\ldots{} up to 30~s before its next claim \\
Anything else & \code{INTERNAL\_ERROR} & Requeued, one attempt more \\
\bottomrule
\end{tabularx}
\caption{Failures of the queue writer.}
\end{table}

If completing or failing a book finds that its task changed, the writer had lost it (the Control requeued it after its
lease expired), and the result is dropped. When the client loses the cluster, the writer stops claiming at once, and
takes the lock again before writing anything when the cluster is back. In real failover runs, a standby replica
took over 2.8~s after a \code{SIGTERM} of the holder and 4.0~s after a \code{kill -9}, and the 17 books the dead holder had
claimed came back through their leases.

\subsection{Idempotency}

Indexing a book twice leaves the index exactly as indexing it once, so every retry is safe:
\begin{itemize}
  \item the posting lists use \code{\$addToSet}, so an id is never added twice, and the unique index on \code{term} keeps
        one document per term;
  \item \code{books} and \code{index\_log} are upserts by \code{\_id};
  \item two replicas, or two threads, writing the same new term at once meet on the unique index. MongoDB retries such
        an upsert by itself because its filter is an equality on the unique key, so neither fails.
\end{itemize}
A test with two \code{IndexingService} of 4 threads each, indexing the same books at once against a real MongoDB,
checks that the index ends as after a single run. \code{docs/robustness.md} lists each failure, what must happen and
the test that proves it, and the mutants those tests catch.

\subsection{Publication of completed index updates}

Search must never show a book half written, and must see all the books of a version at once. A version is a document
\code{CURRENT} in \code{index\_meta}, and a book is visible when \code{index\_books} gives it a version $\le$
\code{CURRENT}. \code{Publisher} publishes on top of \code{CURRENT} = V in four steps:
\begin{enumerate}
  \item the books to publish are those of \code{index\_log} that are not in \code{index\_books}: complete by
        construction;
  \item it records them in \code{index\_books} with version V+1, with \code{\$setOnInsert}: a book keeps forever the
        first version in which it became visible. Search still shows V;
  \item it switches \code{CURRENT} to V+1 \emph{only if it is still V}, a conditional update. This is the instant the
        books become visible, all at once;
  \item it writes the entry of V+1 in \code{index\_publications}, the history.
\end{enumerate}
A process that dies between steps 2 and 3 leaves books above \code{CURRENT}, invisible, and the next publication
completes them. If the conditional switch fails, another replica published first, and the publication is repeated on
top of the new \code{CURRENT}, up to 3 times. The Control decides \emph{when} to publish
(Section~\ref{sec:ingestion-control}); Indexing does it.

After answering, the publication is announced from a thread of its own: the \code{IndexPublished} event on
\code{index-events}, one per version and in order while connected, and the transition of the published books from
\code{INDEXED} to \code{PUBLISHED} in Hazelcast. A slow Hazelcast therefore never holds the publication nor its HTTP
request. Without the cluster the last version is kept and announced on reconnection; a lost event does no harm,
because Search also checks \code{CURRENT} every 5~s.

\subsection{Full rebuild of the index}

\code{POST /index/rebuild} indexes every complete book of the datalake again into a \emph{new} collection,
\code{inverted\_index\_rNN}, while Search keeps reading the current one. It is needed when the tokenizer changes, for
instance. Any replica can start it:
\begin{enumerate}
  \item it puts a marker in Hazelcast (\code{putIfAbsent} with a TTL of 30~s), so only one rebuild runs. While the
        marker exists no writer claims from the queue, and update and publish answer 409 in every replica. A thread of
        the rebuild renews the TTL every 10~s, also in the middle of a batch; if the replica dies, the marker expires
        and indexing goes on;
  \item it takes \code{index-writer}, which the queue writer releases at the end of its batch, and indexes every book
        into \code{rNN} in batches;
  \item a safety pass indexes into \code{rNN} every book published meanwhile in the old collection (a replica without
        Hazelcast does not see the marker);
  \item it switches with one more publication, V+1 on \code{rNN} with every book: Search sees the new collection, and
        every book, at once;
  \item the old collection is dropped after \code{INDEX\_REBUILD\_DROP\_SECONDS} (60~s), well above the 5~s in which a
        Search replica still serves its old snapshot.
\end{enumerate}
If anything fails before the switch, \code{rNN} is dropped and \code{CURRENT} is untouched.

\textbf{Availability.} \code{GET /status} never touches MongoDB, so it answers from the first second. The connection to
Hazelcast is made in the background after the HTTP server starts; until then, and whenever the cluster is lost,
\code{POST /index/update} and \code{POST /index/publish} keep working, because indexing a book does not need the
cluster. If MongoDB does not answer, they answer 503 \code{STORAGE\_UNAVAILABLE}, and retrying is safe.
